\documentclass[11pt]{article}

\usepackage[utf8]{inputenc}
\usepackage[T1]{fontenc}
\usepackage{amsmath,amssymb,amsthm,mathtools}
\usepackage{stmaryrd}
\usepackage[margin=1in]{geometry}
\usepackage{enumitem}
\usepackage{xcolor}
\usepackage[colorlinks=true,linkcolor=blue!60!black,citecolor=blue!60!black,urlcolor=blue!60!black]{hyperref}
\usepackage[nameinlink,capitalize]{cleveref}

\sloppy
\setlength{\emergencystretch}{3em}

\newtheorem{theorem}{Theorem}
\newtheorem{lemma}[theorem]{Lemma}
\newtheorem{corollary}[theorem]{Corollary}
\theoremstyle{definition}
\newtheorem{definition}[theorem]{Definition}
\newtheorem{remark}[theorem]{Remark}

\newcommand{\A}{\mathcal A}
\newcommand{\K}{\mathcal K}
\newcommand{\R}{\mathbb R}
\newcommand{\Homp}{\operatorname{Hom}^{+}}
\newcommand{\Ext}{\operatorname{Ext}}
\newcommand{\brak}[1]{\left\langle #1\right\rangle}
\newcommand{\avg}[2]{\left\llbracket #1 \right\rrbracket_{#2}}

\title{A Two-Root Obstruction for the \texorpdfstring{$C_5$}{C5}-Free Class}
\author{}
\date{\today}

\begin{document}
\maketitle

\begin{abstract}
The one-root planting construction for $C_5$-free graphs does not extend to all
types.  At the two-root edge type there is already a pinning obstruction.  In
every positive-edge-density $C_5$-free limit, a uniformly random rooted edge has
asymptotically zero common-neighbour density, because $C_5$-free graphs have
only $O(e(G))$ triangles.  But the quotient space contains the rooted limit of
a book graph: a distinguished edge with linearly many common neighbours.  Thus
the edge-rooted triangle flag is pinned to $0$ on ensemble supports and attains
$1$ in the quotient, so the $C_5$-free class is not root-plantable at the edge
type.
\end{abstract}

\section{Setup}

Let
\[
  \K=\{G:\text{$G$ contains no $C_5$ as a subgraph}\}.
\]
Let $\tau$ be the two-vertex edge type.  Write $g\in\A^\tau$ for the
three-vertex $\tau$-flag in which the one unlabelled vertex is adjacent to both
labelled vertices; equivalently, $g$ is the triangle flag rooted at an edge.

\section{Triangle density in \texorpdfstring{$C_5$}{C5}-free graphs}

\begin{lemma}[Few triangles]\label{lem:few-triangles}
If $G$ is $C_5$-free, then the number $T(G)$ of triangles in $G$ satisfies
\[
  T(G)\le \frac23 e(G).
\]
In particular, every unlabelled $C_5$-free limit has triangle density $0$.
\end{lemma}

\begin{proof}
For every vertex $v$, the neighbourhood graph $G[N(v)]$ contains no copy of
$P_4$ as a subgraph: a path $a-b-c-d$ inside $N(v)$ would give the cycle
$v-a-b-c-d-v$.  Therefore each connected component of $G[N(v)]$ is either a
star or a triangle, and hence
\[
  e(G[N(v)])\le |N(v)|=d(v).
\]
Summing over $v$ gives
\[
  3T(G)=\sum_{v\in V(G)}e(G[N(v)])
  \le \sum_{v\in V(G)}d(v)=2e(G).
\]
Since $e(G)\le \binom{|V(G)|}{2}$, the normalized triangle density is
$O(1/|V(G)|)$ along any sequence with $|V(G)|\to\infty$.
\end{proof}

\begin{corollary}[The rooted triangle is pinned on ensembles]\label{cor:pinned}
Let $\phi_0\in Q_0$ be a $C_5$-free unlabelled homomorphism with
$\phi_0(\brak{\tau})>0$, and let $\phi^\tau\sim\Ext_\tau(\phi_0)$.  Then
\[
  \phi^\tau(g)=0
  \qquad\text{almost surely.}
\]
\end{corollary}

\begin{proof}
The downward average $\avg{g}{\tau}$ is a positive scalar multiple of the
unlabelled triangle.  By \Cref{lem:few-triangles},
$\phi_0(\avg{g}{\tau})=0$.  Razborov's extension identity gives
\[
  \mathbb E[\phi^\tau(g)]
  =
  \frac{\phi_0(\avg{g}{\tau})}{\phi_0(\brak{\tau})}
  =0.
\]
Since $\phi^\tau(g)\ge0$ pointwise, it is almost surely $0$.
\end{proof}

\section{The quotient sees a book edge}

\begin{definition}[Book graphs]
Let $B_n$ be the graph with two distinguished vertices $a,b$, the edge $ab$,
and $n$ further vertices each adjacent to both $a$ and $b$, with no edges among
the further vertices.
\end{definition}

\begin{lemma}\label{lem:book-c5-free}
Each $B_n$ is $C_5$-free.  Rooted at the ordered edge $(a,b)$, the sequence
$(B_n,a,b)$ has a convergent subsequence whose limit is a point $\psi\in Q_\tau$
with
\[
  \psi(g)=1.
\]
\end{lemma}

\begin{proof}
A leaf of $B_n$ has neighbourhood exactly $\{a,b\}$.  A simple cycle containing
a leaf must enter and leave it through $a$ and $b$.  Hence a simple cycle in
$B_n$ uses at most two leaves, so its length is at most $4$.  Thus $B_n$ has no
$C_5$.

For the rooted density, every unlabelled vertex of $B_n$ is adjacent to both
roots.  Therefore $p(g,(B_n,a,b))=1$ for all $n$, and any convergent subsequence
has the asserted limit.
\end{proof}

\begin{theorem}\label{thm:c5-edge-not-root-plantable}
The $C_5$-free graph class is not root-plantable at the two-root edge type
$\tau$.  Concretely, $-g$ satisfies ensemble semantics but not quotient
semantics.
\end{theorem}

\begin{proof}
By \Cref{cor:pinned}, every support point of every admissible edge-rooted
ensemble has $g=0$, so $-g$ is ensemble-nonnegative.  By
\Cref{lem:book-c5-free}, there is a quotient homomorphism
$\psi\in Q_\tau$ with $\psi(g)=1$, hence $\psi(-g)=-1<0$.  Thus quotient
semantics fails for $-g$, and the support-closure criterion gives
$S_\tau\ne Q_\tau$.
\end{proof}

\begin{remark}
This obstruction is different from the one-root situation.  A single typical
vertex in a positive-density $C_5$-free limit may have nontrivial edge density
to the rest of the graph, and the one-root planting note shows that every such
rooted quotient limit can be planted.  A typical rooted edge is more rigid:
the common-neighbour density is forced to vanish, while an exceptional edge in
a sparse book graph may have common-neighbour density one.
\end{remark}

\end{document}
